\documentclass[12pt]{article}
\usepackage[T1]{fontenc}
\usepackage[utf8]{inputenc}
\usepackage{lmodern}
\usepackage{textcomp}
\usepackage{enumitem}
\usepackage{geometry}
\geometry{margin=1in}
\begin{document}
\begin{center}
Answer Key - Biology - Graduate Level Assessment 6 \
\small (Answers are ordered exactly as in the question paper.)
\end{center}

\section*{Multiple Choice}
\begin{enumerate}[label=\arabic*.]
\item \textbf{Answer:} E
\\ \textit{Options:} A) Through ultrasound examination; B) Through blood sample analysis; C) Through direct observation; D) Through measuring brain wave activity in response to sounds; E) Through operant conditioning
\\ \textit{Note:} Correct option: E - Through operant conditioning
\item \textbf{Answer:} A
\\ \textit{Options:} A) The controls are nerve impulses in the body; B) The controls are the contraction and relaxation cycles of skeletal muscles; C) The controls are hormonal responses that regulate digestion and absorption of food; D) The controls are mechanisms that regulate the heart rate; E) The controls are processes that control the body temperature
\\ \textit{Note:} Correct option: A - The controls are nerve impulses in the body
\item \textbf{Answer:} C
\\ \textit{Options:} A) Appetitive behavior, exploratory behavior, quiescence; B) Termination, appetitive behavior, exploratory behavior; C) Appetitive behavior, consummatory act, quiescence; D) Consummatory act, exploratory behavior, termination; E) Appetitive behavior, termination, consummatory act
\\ \textit{Note:} Correct option: C - Appetitive behavior, consummatory act, quiescence
\item \textbf{Answer:} C
\\ \textit{Options:} A) opt for zero population control once the K value of the curve has been reached; B) maintain the population at the highest point of its logistic curve; C) reduce the carrying capacity cif the environment to lower the K value; D) decrease the mortality rate; E) increase the birth rate of the species
\\ \textit{Note:} Correct option: C - reduce the carrying capacity cif the environment to lower the K value
\item \textbf{Answer:} D
\\ \textit{Options:} A) Coevolution; B) Punctuated equilibrium; C) Directional selection; D) Divergent evolution; E) Convergent evolution
\\ \textit{Note:} Correct option: D - Divergent evolution
\end{enumerate}

\section*{True / False}
\begin{enumerate}[label=\arabic*.]
\setcounter{enumi}{5}
\item \textbf{Answer:} True
\\ \textit{Options:} A) True; B) False
\\ \textit{Note:} The short- and long-term outcomes after LRT and SLT did not differ significantly. To avoid the risk for the donor in LRT, SLT represents the first-line therapy in pediatric liver transplantation in countries where cadaveric organs are available. LRT provides a solution for urgent cases in which a cadaveric graft cannot be found in time or if the choice of the optimal time point for transplantation is vital.
\item \textbf{Answer:} True
\\ \textit{Options:} A) True; B) False
\\ \textit{Note:} The available evidence suggests that the operative treatment for DMCF is associated with a lower rate of nonunion, malunion and complication than nonoperative treatment. This study supports traditional primary operative treatment for DMCF in active adults.
\item \textbf{Answer:} True
\\ \textit{Options:} A) True; B) False
\\ \textit{Note:} For some patients, using colour to describe their pain experience may be a useful tool to improve doctor-patient communication.
\item \textbf{Answer:} True
\\ \textit{Options:} A) True; B) False
\\ \textit{Note:} Elective re-siting of intravenous cannulae every 48 hours results in a significant reduction in the incidence and severity of PVT. We recommend that this should be adopted as standard practice in managing all patients who require prolonged intravenous therapy.
\item \textbf{Answer:} True
\\ \textit{Options:} A) True; B) False
\\ \textit{Note:} When there is clinical suspicion of sepsis, appropriate empirical systemic antibiotic therapy should be broad spectrum and should rely on the susceptibility of the organisms from recent cultures of the burn wound surface, until the blood cultures results are completed.
\end{enumerate}

\section*{Long Form Response}
\begin{enumerate}[label=\arabic*.]
\setcounter{enumi}{10}
\item \textbf{Answer:} Cortisol, aldosterone, and adrenal androgens. Explain why this is the correct answer and provide supporting evidence. Reference key facts or concepts that support this choice.
\\ \textit{Note:} Long-form derived from MCQ; award credit for stating the correct idea and giving supporting reasoning.
\item \textbf{Answer:} mRNA would replicate rapidly. Explain why this is the correct answer and provide supporting evidence. Reference key facts or concepts that support this choice.
\\ \textit{Note:} Long-form derived from MCQ; award credit for stating the correct idea and giving supporting reasoning.
\end{enumerate}

\vfill
\noindent\textit{End of Answer Key}
\end{document}